	\section{Proof of the main result}
	
	In this section, we show how to deduce Theorem \ref{thm:main0} from Theorem \ref{thm:main}. Let us start with the following variant of the well known Oddtown theorem \cite{B69}, see also \cite{BF92} for related results.
	
	\begin{lemma}\label{lemma:oddtown}
		Let $\ell,m,n$ be positive integers, and let $A_1,\dots,A_m,B_1,\dots,B_m\subset [n]$ such that $\ell\nmid|A_{i}\cap B_{i}|$ for $i\in [m]$, but $\ell$ divides $|A_{i}\cap B_{j}|$ for $i\neq j$. Then $m\leq sn$, where $s$ is the number of distinct prime divisors of~$\ell$.
	\end{lemma}
	
	\begin{proof}
		Write $\ell=p_1^{\alpha_1}\dots p_{s}^{\alpha_s}$, where $p_1,\dots,p_s$ are distinct primes. Let $v_{i}$ and $w_{i}$ be the characteristic vectors of $A_{i}$ and $B_{i}$ over $\mathbb{Q}$, respectively. Let $t=\lceil m/s\rceil$, then there exists $r\in [s]$ such that for at least $t$ of the indices $i\in [m]$, we have that $|A_{i}\cap B_{i}|$ is not divisible by $p_{r}^{\alpha_r}$. Without loss of generality let these $t$ indices be $1,\dots,t$. We show that $v_1,\dots,v_t$ are linearly independent (over $\mathbb{Q}$), which then implies $t\leq n$ and $m\leq sn$.
		
		Suppose this is not the case, then there exist $c_1,\dots,c_t\in\mathbb{Z}$, not all zero, such that $\sum_{i=1}^{t}c_{i}v_{i}=0$. We can assume that at least one of $c_1,\dots,c_t$ is not divisible by $p_r$, otherwise we can replace $c_i$ with $c_i'=c_i/p_r$ for every $i\in [t]$. Let $k\in [t]$ be an index such that $p_r\nmid c_{k}$. Consider the equality 
		$$0=\left\langle\sum_{i=1}^{t}c_{i} v_{i},w_{k}\right\rangle=\sum_{i=1}^{t}c_{i} |A_i\cap B_k|.$$
		We have $p_r^{\alpha_{r}}\mid c_i|A_{i}\cap B_{k}|$ if $i\neq k$, and $p_r^{\alpha_r}\nmid c_k|A_{k}\cap B_{k}|$, so $p_r^{\alpha_{r}}\nmid \langle\sum_{i=1}^{t}c_{i} v_{i},w_{k}\rangle$, contradiction.
	\end{proof}
	
	Say that $\mathcal{F}\subset 2^{[n]}$ is \emph{weakly $k$-closed over $\mathbb{Z}_{\ell}$} if the intersection of any $k$ distinct elements of $\mathcal{F}$ is divisible by $\ell$. Also, say that $\mathcal{F}\subset 2^{[n]}$ is \emph{$k$-closed over $\mathbb{Z}_{\ell}$} if the family formed by the characteristic vectors of the elements of $\mathcal{F}$ is $k$-closed over $\mathbb{Z}_{\ell}$. So $\mathcal{F}$ is $k$-closed over $\mathbb{Z}_{\ell}$ if and only if the intersection of any $k$ not necessarily distinct elements of $\mathcal{F}$ is divisible by $\ell$. We need the following useful observation, which for $\ell=2$ appears in \cite{SV18}. 
	
	\begin{lemma}\label{lemma:weakly}
		Let $\ell,k$ be positive integers, and let $s$ be the number of distinct prime divisors of $\ell$. Let $\mathcal{F}\subset 2^{[n]}$ such that $\mathcal{F}$ is weakly $k$-closed over $\mathbb{Z}_{\ell}$.  Then there exists $\mathcal{F}'\subset \mathcal{F}$ such that $|\mathcal{F}'|\geq |\mathcal{F}|-sk^2n$, and $\mathcal{F}'$ is $k$-closed over $\mathbb{Z}_{\ell}$. 
	\end{lemma}
	
	\begin{proof}
		Repeat the following removal operation. Suppose that $\mathcal{F}$ is not $k$-closed, and let $t$ be maximal such that some $t$ distinct elements of $\mathcal{F}$ have an intersection not divisible by $\ell$. So $t < k$ because $\mathcal{F}$ is weakly $k$-closed. Let $\mathcal{H}$ be the $t$-uniform hypergraph on $\mathcal{F}$ in which $\{C_{1},\dots,C_{t}\}$ is an edge if $|C_1\cap\dots\cap C_t|$ is not divisible by $\ell$. We claim that $\mathcal{H}$ contains no matching of size more than $sn$. Indeed, suppose otherwise, let $\{C_{i,1},\dots,C_{i,t}\}$, $i\in [m]$, be the edges of a matching of size $m>sn$. For $i\in [m]$, let $A_{i}=C_{i,1}$ and $B_{i}=C_{i,1}\cap\dots\cap C_{i,t}$. Then $\ell\nmid|A_{i}\cap B_{i}|$, but $\ell$ divides $|A_i\cap B_j|$ for every $i\neq j$ by the maximality of $t$. Therefore, by Lemma \ref{lemma:oddtown} we get $m\leq sn$, contradiction.
		
		Consider a maximal matching of $\mathcal{H}$, and remove every element of $\mathcal{F}$ that appears in this matching. Then $\mathcal{F}$ no longer contains $t$ distinct sets, whose intersection is not divisible by $\ell$. Repeating this procedure at most $k-1$ times, we get a family $\mathcal{F}'\subset \mathcal{F}$ such that $\mathcal{F}'$ is $k$-closed over $\mathbb{Z}_{\ell}$, and $|\mathcal{F}'|\geq |\mathcal{F}|-sk^2n$.
	\end{proof}
	
	Lemma \ref{lemma:weakly} combined with Theorem \ref{thm:main} immediately implies that if $\mathcal{F}\subset 2^{[n]}$ is weakly $k$-closed over $\mathbb{Z}_{\ell}$, then $|\mathcal{F}|\leq 2^{\lfloor n/\ell\rfloor}+sk^2n$. In order to improve the term $sk^2n$ to a constant, we use the second part of Theorem~\ref{thm:main}.
	
	\begin{proof}[Proof of Theorem \ref{thm:main0}]
		Let $d=\lfloor n/\ell\rfloor$. Let $\mathcal{F}\subset\{0,1\}^{n}$ such that $\mathcal{F}$ is weakly $k$-closed over $\mathbb{Z}_{\ell}$. Then by Lemma \ref{lemma:weakly}, there exists $\mathcal{F}'\subset \mathcal{F}$ such that $|\mathcal{F}'|\geq |\mathcal{F}|-sk^2n$ and $\mathcal{F}'$ is $k$-closed over $\mathbb{Z}_{\ell}$, where $s$ is the number of distinct prime divisors of $\ell$. If $|\mathcal{F}'|\leq 2^{d-1}$, then 
		$$|\mathcal{F}|\leq 2^{d-1}+sk^2n<2^{d},$$
		where the last inequality holds if $n$ is sufficiently large. 
		
		Suppose that $|\mathcal{F}'|> 2^{d-1}$ and $|\mathcal{F}|\geq 2^{d}$, otherwise we are done. Then by Theorem \ref{thm:main}, $[n]$ can be partitioned into sets $A_{1},\dots,A_{d},A'$ such that $A_{i}$ is a maximal set of twins for $\mathcal{F}'$ for $i\in [d]$, $|A_{i}|=\ell$, $|A'|\leq \ell-1$, and $\mathcal{F}'$ vanishes on $A'$. Let $S\subset \{0,1\}^{n}$ be the atomic family containing all possible $2^{d}$ sets $C$ such that $C\cap A_{i}\in \{\emptyset,A_{i}\}$ for every $i\in [d]$. Then $\mathcal{F}'\subset S$ and $|S\setminus\mathcal{F}'|\leq sk^2n$, as $|\mathcal{F}'| \geq 2^d - sk^2n$. Also, if $n$ is sufficiently large, for every $i\in [d]$ we can find $k-1$ distinct sets $B_{i,1},\dots,B_{i,k-1}\in \mathcal{F}'$ such that $A_i=\bigcap_{j=1}^{k-1}B_{i,j}$. Indeed, $\mathcal{F}'$ contains a set of the form $A_{i}\cup A_{a}\cup A_{b}$ for some $a,b$, as the number of such sets in $S$ is $\binom{d-1}{2}>sk^2n$. Let $B_{i,1}$ be such a set $A_{i}\cup A_{a}\cup A_{b} \in \mathcal{F}'$. Also $\mathcal{F}'$ contains $k-2$ sets that contain $A_i$ but do not contain $A_{a}$ and $A_b$ as the number of such sets in $S$ is $2^{d-3}> sk^2n+k$. Let these $k-2$ sets be $B_{i,2},\dots,B_{i,k-1}$. Then $A_{i}=\bigcap_{j=1}^{k-1}B_{i,j}$, as claimed. 
		
		Consider any $F\in \mathcal{F}\setminus S$.
		For every $i\in [d]$, we have $A_i\subset F$ or $A_i\cap F=\emptyset$, as the size of  $A_i\cap F=B_{i,1}\cap\dots\cap B_{i,k-1}\cap F$ must be divisible by $\ell$. 
		Now, as $F \notin S$, we must have 
		$F \cap A' \neq \emptyset$.
		On the other hand, for any $H\subset A'$ with $H\neq \emptyset$, there are at most $k-1$ elements $F\in\mathcal{F}\setminus S$ such that $F\cap A'=H$, because otherwise we would have $k$ distinct $F_1,\dots,F_k \in \mathcal{F}$ with $|F_1 \cap \dots \cap F_k| \equiv |H| \not\equiv 0 \pmod{\ell}$, a contradiction. 
		So we see that $|\mathcal{F}\setminus S|\leq k2^{|A'|}\leq k2^{\ell-1}$, and if $\ell \mid n$ then $\mathcal{F}\subset S$ because $A' = \emptyset$. This finishes the proof.
	\end{proof}
